% GENERATED FILE. Edit canonical lessons, then run npm run generate:books.
% canonical-source-hash: fnv1a64:d617281ea8d9bba0
% canonical-lessons: KA-C367-khalanayaka, KA-C367-vira, KA-C367-satru, KA-C367-nasamadu, KA-C367-kanmareyagu

\chapter{Heroes and Villains}
\label{ch:ka-heroes-and-villains}
\hlchaptermodality{367}

\begin{chapteropening}
\textbf{By the end of this chapter:} \emph{I can say a villain, a hero, an enemy, to destroy and to disappear.}

Complete the last lesson of chapter 367: I can say a villain, a hero, an enemy, to destroy and to disappear.
\end{chapteropening}

\section[\texorpdfstring{\kn{ಖಳನಾಯಕ} (khaḷanāyaka)}{ಖಳನಾಯಕ (khaḷanāyaka)}]{\kn{ಖಳನಾಯಕ} (khaḷanāyaka) — a villain}

\label{lesson:KA-C367-khalanayaka}



\begin{quote}
Before the new one: say the Kannada for an oath, then the Kannada for to wail.
\end{quote}

\subsection*{You'll want to know: \kn{ಖಳನಾಯಕ}}
\textbf{\kn{ಖಳನಾಯಕ}} — \emph{khaḷanāyaka} — \textquotedblleft{}a villain\textquotedblright{}.

\begin{quote}
\kn{ಸಿನಿಮಾದ} \kn{ಕೊನೆಯಲ್ಲಿ} \kn{ಖಳನಾಯಕ} \kn{ಸೋತನು}.
\end{quote}

\begin{quote}
\emph{sinimāda koneyalli khaḷanāyaka sōtanu.}
\end{quote}

\begin{quote}
At the end of the film the villain lost.
\end{quote}

Plural: \kn{ಖಳನಾಯಕರು}, \emph{khaḷanāyakaru}.

\kn{ಖಳ} is wicked: the wicked lead of a story or a film.

\begin{cousinweb}
\textbf{\kn{ಖಳನಾಯಕ}} (\emph{khaḷanāyaka}) is borrowed from Sanskrit, a compound of \emph{khala}, \textquotedblleft{}a mischievous person\textquotedblright{} (CDIAL 3835), and \emph{nāyaka}, \textquotedblleft{}a guide, a leader\textquotedblright{} (Monier-Williams).
\end{cousinweb}

\begin{grammarlens}[title={What you've built}]
One more word for telling what happened.
\end{grammarlens}

\subsection*{Your turn}
\begin{itemize}
\raggedright
  \item \emph{Say it:} \emph{khaḷanāyaka}
  \item \emph{Say it:} \emph{khaḷanāyaka}, once more
  \item \emph{Say it:} \emph{sinimāda koneyalli khaḷanāyaka sōtanu.} — At the end of the film the villain lost.
  \item \emph{Recall:} say the Kannada for an oath, then the Kannada for to wail, then say \emph{khaḷanāyaka} again
\end{itemize}

\begin{tcolorbox}[breakable,colback=teal!4,colframe=teal!35!black,title={Before you move on}]
What does \kn{ಖಳನಾಯಕ} mean? (\textquotedblleft{}A villain\textquotedblright{}.) Say it once more.
\end{tcolorbox}
\section[\texorpdfstring{\kn{ವೀರ} (vīra)}{ವೀರ (vīra)}]{\kn{ವೀರ} (vīra) — a hero; a brave man}

\label{lesson:KA-C367-vira}



\begin{quote}
Before the new one: say the Kannada for to wail, then the Kannada for a villain.
\end{quote}

\subsection*{You'll want to know: \kn{ವೀರ}}
\textbf{\kn{ವೀರ}} — \emph{vīra} — \textquotedblleft{}a hero; a brave man\textquotedblright{}.

\begin{quote}
\kn{ಆ} \kn{ವೀರ} \kn{ಒಬ್ಬನೇ} \kn{ಹಳ್ಳಿಯನ್ನು} \kn{ಕಾಪಾಡಿದನು}.
\end{quote}

\begin{quote}
\emph{ā vīra obbanē haḷḷiyannu kāpāḍidanu.}
\end{quote}

\begin{quote}
That hero saved the village on his own.
\end{quote}

Plural: \kn{ವೀರರು}, \emph{vīraru}.

As an adjective: \kn{ವೀರ} \kn{ಯೋಧ}, a brave warrior. A heroine is a \kn{ವೀರೆ} or \kn{ವೀರ} \kn{ಮಹಿಳೆ}.

\begin{cousinweb}
\textbf{\kn{ವೀರ}} (\emph{vīra}) is borrowed from Sanskrit \emph{vīra}, \textquotedblleft{}a man, a hero\textquotedblright{} (CDIAL 12056).
\end{cousinweb}

\begin{grammarlens}[title={What you've built}]
One more word for telling what happened.
\end{grammarlens}

\subsection*{Your turn}
\begin{itemize}
\raggedright
  \item \emph{Say it:} \emph{vīra}
  \item \emph{Say it:} \emph{vīra}, once more
  \item \emph{Say it:} \emph{ā vīra obbanē haḷḷiyannu kāpāḍidanu.} — That hero saved the village on his own.
  \item \emph{Recall:} say the Kannada for to wail, then the Kannada for a villain, then say \emph{vīra} again
\end{itemize}

\begin{tcolorbox}[breakable,colback=teal!4,colframe=teal!35!black,title={Before you move on}]
What does \kn{ವೀರ} mean? (\textquotedblleft{}A hero; a brave man\textquotedblright{}.) Say it once more.
\end{tcolorbox}
\section[\texorpdfstring{\kn{ಶತ್ರು} (śatru)}{ಶತ್ರು (śatru)}]{\kn{ಶತ್ರು} (śatru) — an enemy}

\label{lesson:KA-C367-satru}



\begin{quote}
Before the new one: say the Kannada for a villain, then the Kannada for a hero.
\end{quote}

\subsection*{You'll want to know: \kn{ಶತ್ರು}}
\textbf{\kn{ಶತ್ರು}} — \emph{śatru} — \textquotedblleft{}an enemy\textquotedblright{}.

\begin{quote}
\kn{ವೀರ} \kn{ತನ್ನ} \kn{ಶತ್ರುವನ್ನು} \kn{ಕ್ಷಮಿಸಿದನು}.
\end{quote}

\begin{quote}
\emph{vīra tanna śatruvannu kṣamisidanu.}
\end{quote}

\begin{quote}
The hero forgave his enemy.
\end{quote}

Plural: \kn{ಶತ್ರುಗಳು}, \emph{śatrugaḷu}.

A friend is a \kn{ಮಿತ್ರ}; \kn{ಶತ್ರು} is the opposite.

\begin{cousinweb}
\textbf{\kn{ಶತ್ರು}} (\emph{śatru}) is borrowed from Sanskrit \emph{śatru}, \textquotedblleft{}an enemy\textquotedblright{} (CDIAL 12286).
\end{cousinweb}

\begin{grammarlens}[title={What you've built}]
One more word for telling what happened.
\end{grammarlens}

\subsection*{Your turn}
\begin{itemize}
\raggedright
  \item \emph{Say it:} \emph{śatru}
  \item \emph{Say it:} \emph{śatru}, once more
  \item \emph{Say it:} \emph{vīra tanna śatruvannu kṣamisidanu.} — The hero forgave his enemy.
  \item \emph{Recall:} say the Kannada for a villain, then the Kannada for a hero, then say \emph{śatru} again
\end{itemize}

\begin{tcolorbox}[breakable,colback=teal!4,colframe=teal!35!black,title={Before you move on}]
What does \kn{ಶತ್ರು} mean? (\textquotedblleft{}An enemy\textquotedblright{}.) Say it once more.
\end{tcolorbox}
\section[\texorpdfstring{\kn{ನಾಶಮಾಡು} (nāśamāḍu)}{ನಾಶಮಾಡು (nāśamāḍu)}]{\kn{ನಾಶಮಾಡು} (nāśamāḍu) — to destroy}

\label{lesson:KA-C367-nasamadu}



\begin{quote}
Before the new one: say the Kannada for a hero, then the Kannada for an enemy.
\end{quote}

\subsection*{You'll want to know: \kn{ನಾಶಮಾಡು}}
\textbf{\kn{ನಾಶಮಾಡು}} — \emph{nāśamāḍu} — \textquotedblleft{}to destroy\textquotedblright{}.

\begin{quote}
\kn{ಖಳನಾಯಕ} \kn{ಇಡೀ} \kn{ಊರನ್ನು} \kn{ನಾಶಮಾಡಲು} \kn{ಬಯಸಿದನು}.
\end{quote}

\begin{quote}
\emph{khaḷanāyaka iḍī ūrannu nāśamāḍalu bayasidanu.}
\end{quote}

\begin{quote}
The villain wanted to destroy the whole town.
\end{quote}

Past (I …): \kn{ನಾಶಮಾಡಿದೆ}, \emph{nāśamāḍide}.

\kn{ನಾಶ}, ruin, with \kn{ಮಾಡು}. Something destroyed \kn{ನಾಶವಾಯಿತು}.

\begin{cousinweb}
\textbf{\kn{ನಾಶಮಾಡು}} (\emph{nāśamāḍu}) is built from Sanskrit \emph{nāśa}, \textquotedblleft{}destruction, ruin\textquotedblright{} (CDIAL 7084), and native \emph{māḍu}, \textquotedblleft{}to do\textquotedblright{} (DEDR 4797).
\end{cousinweb}

\begin{grammarlens}[title={What you've built}]
One more word for telling what happened.
\end{grammarlens}

\subsection*{Your turn}
\begin{itemize}
\raggedright
  \item \emph{Say it:} \emph{nāśamāḍu}
  \item \emph{Say it:} \emph{nāśamāḍu}, once more
  \item \emph{Say it:} \emph{khaḷanāyaka iḍī ūrannu nāśamāḍalu bayasidanu.} — The villain wanted to destroy the whole town.
  \item \emph{Recall:} say the Kannada for a hero, then the Kannada for an enemy, then say \emph{nāśamāḍu} again
\end{itemize}

\begin{tcolorbox}[breakable,colback=teal!4,colframe=teal!35!black,title={Before you move on}]
What does \kn{ನಾಶಮಾಡು} mean? (\textquotedblleft{}To destroy\textquotedblright{}.) Say it once more.
\end{tcolorbox}
\section[\texorpdfstring{\kn{ಕಣ್ಮರೆಯಾಗು} (kaṇmareyāgu)}{ಕಣ್ಮರೆಯಾಗು (kaṇmareyāgu)}]{\kn{ಕಣ್ಮರೆಯಾಗು} (kaṇmareyāgu) — to disappear; to vanish}

\label{lesson:KA-C367-kanmareyagu}



\begin{quote}
Before the new one: say the Kannada for an enemy, then the Kannada for to destroy.
\end{quote}

\subsection*{You'll want to know: \kn{ಕಣ್ಮರೆಯಾಗು}}
\textbf{\kn{ಕಣ್ಮರೆಯಾಗು}} — \emph{kaṇmareyāgu} — \textquotedblleft{}to disappear; to vanish\textquotedblright{}.

\begin{quote}
\kn{ಮಾಂತ್ರಿಕ} \kn{ಹೊಗೆಯಲ್ಲಿ} \kn{ಕಣ್ಮರೆಯಾದನು}.
\end{quote}

\begin{quote}
\emph{māntrika hogeyalli kaṇmareyādanu.}
\end{quote}

\begin{quote}
The sorcerer vanished in the smoke.
\end{quote}

Past (it …): \kn{ಕಣ್ಮರೆಯಾಯಿತು}, \emph{kaṇmareyāyitu}.

\kn{ಕಣ್}, eye, \kn{ಮರೆ}, hidden, and \kn{ಆಗು}: to become hidden from the eye.

\begin{cousinweb}
\textbf{\kn{ಕಣ್ಮರೆಯಾಗು}} (\emph{kaṇmareyāgu}) is a native Dravidian compound: \emph{kaṇ}, \textquotedblleft{}eye\textquotedblright{} (DEDR 1159), \emph{mare}, \textquotedblleft{}hidden, screened\textquotedblright{} (DEDR 4760), and \emph{āgu}, \textquotedblleft{}to become\textquotedblright{} (DEDR 333).
\end{cousinweb}

\begin{grammarlens}[title={What you've built}]
One more word for telling what happened.
\end{grammarlens}

\subsection*{Your turn}
\begin{itemize}
\raggedright
  \item \emph{Say it:} \emph{kaṇmareyāgu}
  \item \emph{Say it:} \emph{kaṇmareyāgu}, once more
  \item \emph{Say it:} \emph{māntrika hogeyalli kaṇmareyādanu.} — The sorcerer vanished in the smoke.
  \item \emph{Recall:} say the Kannada for an enemy, then the Kannada for to destroy, then say \emph{kaṇmareyāgu} again
\end{itemize}

\begin{tcolorbox}[breakable,colback=teal!4,colframe=teal!35!black,title={Before you move on}]
What does \kn{ಕಣ್ಮರೆಯಾಗು} mean? (\textquotedblleft{}To disappear; to vanish\textquotedblright{}.) Say it once more.
\end{tcolorbox}
